% Imported from topics/video-generation/seedance-2.5/integrated-training.md
\chapter{M20 联合模型训练课程}\label{seedance-m20}
\moduleinfo{M20}{候选设计与工程选型；组合效果待验证}{\hyperref[src-ltx-2]{\texttt{ltx-2}}、\hyperref[src-vace]{\texttt{vace}}、\hyperref[src-alchemint]{\texttt{alchemint}}、\hyperref[src-videocof]{\texttt{videocof}}、\hyperref[src-vincie-next]{\texttt{vincie-next}}、\hyperref[src-long-context-tuning]{\texttt{long-context-tuning}}、\hyperref[src-longlive]{\texttt{longlive}}、\hyperref[src-rewarddance]{\texttt{rewarddance}}、\hyperref[src-ltx-original-code]{\texttt{ltx-original-code}}}


参考、编辑、空间控制和长叙事不应各训一个模型后再假定权重可以合并。一个完整候选流程要在同一潜空间和主干上安排任务，再通过课程决定何时增加难度、解冻参数和改变注意力结构。

这里选用 LTX-2 原系列非蒸馏联合模型，保留其视频/音频VAE、文本组件和主干，训练新增参考与源条件支路，再逐步扩大到长片。下述比例与权重都是初始候选配方，未进行模型运行、训练或预算实测。

\section{每种能力需要什么样本}\label{seedance-m20-ux6bcfux79cdux80fdux529bux9700ux8981ux4ec0ux4e48ux6837ux672c}

{\def\LTcaptype{table} % do not increment counter
\small\setlength{\tabcolsep}{3pt}\renewcommand{\arraystretch}{1.18}
\begin{longtable}[]{@{}
  >{\raggedright\arraybackslash}p{(\linewidth - 4\tabcolsep) * \real{0.1900}}
  >{\raggedright\arraybackslash}p{(\linewidth - 4\tabcolsep) * \real{0.3900}}
  >{\raggedright\arraybackslash}p{(\linewidth - 4\tabcolsep) * \real{0.4200}}@{}}
\toprule\noalign{}
\begin{minipage}[b]{\linewidth}\raggedright
数据池
\end{minipage} & \begin{minipage}[b]{\linewidth}\raggedright
最小样本
\end{minipage} & \begin{minipage}[b]{\linewidth}\raggedright
关键监督与排错
\end{minipage} \\
\midrule\noalign{}
\endhead
\bottomrule\noalign{}
\endlastfoot
基础音视频 & caption、视频、音轨、对齐时间 & 不漏事件、不引入原始声画偏移 \\
身份参考 & 同主体跨片段图像、目标动作、可见区间 & 防同帧复制、区分离场与遮挡 \\
动作/相机 & pose/flow/camera与目标视频 & 控制信号不泄漏目标身份 \\
音频参考 & 同speaker不同台词/情绪、转录与目标音轨 & 分离音色、文本内容、背景音乐 \\
区域编辑 & 源、区域、指令、目标、保持规则 & 目标改变与外区保持同时成立 \\
示例编辑 & 源视频、原关键帧、编辑后关键帧 & 外观示例传播且动作不丢 \\
几何控制 & clay/depth/normal、相机、目标材质渲染 & 同一场景单位、遮挡和时间对齐 \\
场景叙事 & global/shot prompt、角色道具状态、连续片段 & 事件顺序、长程身份与道具状态 \\
\end{longtable}
}

所有派生样本继承来源分组，参考与目标也参与去重关系。音色数据按speaker分组，场景数据按原始视频/三维资产家族分组；不能只对裁出的片段随机划分。

\section{八类数据怎样从素材中做出来}\label{seedance-m20-ux516bux7c7bux6570ux636eux600eux6837ux4eceux7d20ux6750ux4e2dux505aux51faux6765}

先完成媒体入库、真实重采样、声画时间对齐、切镜、去重关系和来源划分，再制作任务配对。最后才缓存latent，避免清洗或裁剪已改变、缓存却仍指向旧媒体。

{\def\LTcaptype{table} % do not increment counter
\small\setlength{\tabcolsep}{3pt}\renewcommand{\arraystretch}{1.18}
\begin{longtable}[]{@{}
  >{\raggedright\arraybackslash}p{(\linewidth - 4\tabcolsep) * \real{0.1900}}
  >{\raggedright\arraybackslash}p{(\linewidth - 4\tabcolsep) * \real{0.3900}}
  >{\raggedright\arraybackslash}p{(\linewidth - 4\tabcolsep) * \real{0.4200}}@{}}
\toprule\noalign{}
\begin{minipage}[b]{\linewidth}\raggedright
数据池
\end{minipage} & \begin{minipage}[b]{\linewidth}\raggedright
可执行的制作顺序
\end{minipage} & \begin{minipage}[b]{\linewidth}\raggedright
需要人工确认的关键点
\end{minipage} \\
\midrule\noalign{}
\endhead
\bottomrule\noalign{}
\endlastfoot
基础音视频 & 对齐时间戳→筛除声画错位→镜头/事件描述→ASR与说话人分段 & 台词是否真实可听，事件是否漏标 \\
身份参考 & 跨片段主体匹配→跟踪与可见性→抽不同视角参考→标出多个出场区间 & 是否同一身份，离场与遮挡是否混淆 \\
动作/相机 & 从目标提取pose/flow/相机→移除外观路径→跨身份配对或受控重渲染 & 控制信号是否夹带目标服装/背景 \\
音频参考 & 确认speaker→分离干净语音段→抽取不同台词参考→保留转录与情绪 & 说话人一致，背景音乐是否泄漏 \\
区域编辑 & 同一场景受控重渲染或人工编辑对→逐帧mask→编辑指令→外区比较 & 变化是否只有标签声明的属性 \\
示例编辑 & 选源关键帧→制作编辑示例→配可验证的视频目标 & 背面和遮挡后外观是否一致 \\
几何控制 & 固定资产/相机→导出clay/depth/normal/ID→渲染目标材质→同步裁剪 & 单位、轴、遮挡和帧映射 \\
场景叙事 & 保留连续场景→分镜→全局/局部描述→记录角色/道具状态 & 状态迁移与事件先后是否真实 \\
\end{longtable}
}

从目标中提取动作或深度是有监督控制任务的合法条件；它不能再被当作``没有这些条件也能预测未来''的证据。图像编辑结果若没有可靠视频目标，只能进入图像任务或标为弱监督，不能伪装成真实视频编辑对。

每条样本同时保存内容ID、来源组、split、全部条件资产、时间变换、task类型、音频状态和监督有效位；speaker/几何/多区间标签作为候选扩展，需落到具体数据加载器，现有基础schema不会自动校验这些新语义。先对每池做分层小样本审阅，修正配对和标注规则，再扩大规模。

\section{统一目标与有效位}\label{seedance-m20-ux7edfux4e00ux76eeux6807ux4e0eux6709ux6548ux4f4d}

以教学记号 \(t=0\) 为干净、\(t=1\) 为噪声，目标视频/音频分别构造：

\[
z_t^m=(1-t_m)z_0^m+t_m\epsilon_m,
\qquad u_m^*=\epsilon_m-z_0^m,\quad m\in\{v,a\}.
\]

实现保留原模型的sigma/time映射。已知条件与目标分别携带时间和denoise mask，不能用同一个全局标量把干净源条件调制成未知目标。

联合生成先沿用原生sigma采样和两流噪声约定。保留音轨时音频固定为干净条件，只有视频是生成目标；缺音轨时音频项与无效交互都关闭。若后续引入两流独立噪声采样，需要专门覆盖相应噪声组合并验证，不能只改公式中的两个下标。

每模态先在每条样本内对有效元素平均，再对有监督的样本取平均，最后合成任务损失：

下面的 \(i\) 遍历标量通道元素，\(W_i^m\) 为其有效权重；若实现中先对token向量求平方和，分母还需计入有效通道数。

\[
\mathcal{L}_{\mathrm{flow}}^m=
\frac{\sum_i W_i^m\lVert u_i^m-u_i^{m*}\rVert^2}
{\sum_i W_i^m},\qquad
\mathcal{L}_{\mathrm{joint}}=\mathcal{L}_{\mathrm{flow}}^v+
\lambda_a\mathcal{L}_{\mathrm{flow}}^a.
\]

缺音轨时跳过音频项，而不是给零分母加epsilon假装存在目标。分布式有效样本数不同时使用全局有效计数修正，不能只平均各卡均值。初始 \(\lambda_a=1\)，随后按有效元素标准化后的梯度和独立音质评价调节，不直接按视频/音频token数量配权。

\section{S0：共同接口与基础能力固定}\label{seedance-m20-s0ux5171ux540cux63a5ux53e3ux4e0eux57faux7840ux80fdux529bux56faux5b9a}

先固定原模型、VAE归一化、文本投影、帧网格、音频token率和所有源码版本。准入阶段应确认原版生成、视频重建与声画对齐符合模型预期，再接新支路；这些是后续实施条件，当前配方尚未经实际执行。完整数据闭环需要真实样本、同步时间与可恢复训练记录，格式样例不能充当训练结果。

新增支路正常初始化投影、零门控，不改变原模型初始输出。视频/音频codec和冻结文本编码器不参与首轮更新。

冻结主干权重仍要保留通向新增支路的梯度路径；不能把后续主干整体放在 \codepath{no_grad()} 中。优化器与保存逻辑显式收集新池化器、条件塔、投影和门控，不能只筛选LoRA权重。

所有后续阶段共用一个更新顺序：抽取任务与形状桶，加载一致时间轴的媒体/条件，构造目标噪声与有效位，执行原主干和新增支路，按任务合成损失，再做梯度累积、裁剪与更新。条件编码器冻结时可以离线缓存；可学习池化器和adapter的输出必须在线计算，不能缓存后继续声称它们在学习。

\section{S1：参考、动作和结构支路}\label{seedance-m20-s1ux53c2ux8003ux52a8ux4f5cux548cux7ed3ux6784ux652fux8def}

训练参考池化器、绑定embedding、参考attention、视觉context adapter和结构投影，主干先冻结。初始有效global batch设为8，adapter学习率 \(10^{-4}\)，AdamW，梯度裁剪1.0；是否启用主干LoRA由固定对照决定，不能把``冻结不收敛''一律解释为数据不足。

{\def\LTcaptype{table} % do not increment counter
\small\setlength{\tabcolsep}{3pt}\renewcommand{\arraystretch}{1.18}
\begin{longtable}[]{@{}
  >{\raggedright\arraybackslash}p{(\linewidth - 2\tabcolsep) * \real{0.2200}}
  >{\raggedleft\arraybackslash}p{(\linewidth - 2\tabcolsep) * \real{0.7800}}@{}}
\toprule\noalign{}
\begin{minipage}[b]{\linewidth}\raggedright
任务
\end{minipage} & \begin{minipage}[b]{\linewidth}\raggedleft
初始采样比例
\end{minipage} \\
\midrule\noalign{}
\endhead
\bottomrule\noalign{}
\endlastfoot
基础T2AV回归 & 20\% \\
图片身份/外观参考 & 30\% \\
视频动作/相机控制 & 20\% \\
音色/声音参考 & 20\% \\
白模与几何控制 & 10\% \\
\end{longtable}
}

这些比例按样本任务采样，再做token预算约束；不能让长样本因计算更多而隐式占据全部更新。若没有已授权音色配对数据，不伪造这一池，改为独立后期配音，并记录该能力未解锁。

从1--2份参考、短片起步，逐步到4/8份，再到设计的30/10/10文件上限。素材数量、有效条件数和压缩token数一起记录。区间外缺席、相似主体重叠、动作参考换身份必须出现于训练和验证。

先安排1000次优化器更新的短pilot、100次预热，每100次固定验证，作为排查接口的候选起点；这些数字不代表足够训练量。BF16在硬件支持时启用，单卡microbatch与累积数按全局batch8配平。首次解冻主干的候选LoRA取rank16、alpha16、学习率 \(10^{-5}\)，仅选原生self与AV注意力的 \codepath{to_q}、\codepath{to_v}，登记完整模块名并排除文本与新增支路。该数值是起点，不代表适合所有权重。

\section{S2：定位与局部编辑}\label{seedance-m20-s2ux5b9aux4f4dux4e0eux5c40ux90e8ux7f16ux8f91}

固定已可用的参考模块，先训练源条件与区域路径；必要时再联合微调适配器和有限主干LoRA。初始混合为：基础T2AV 15\%、参考任务20\%、真实/受控视频编辑30\%、关键帧示例传播20\%、区域定位15\%。

人工mask优先作为受控条件。无mask样本先预测区域，再给第二阶段读取；逐步用预测区域替代部分真值区域，使第二阶段面对真实错误分布。若采用分割头，区域损失可用BCE与Dice；若采用灰度区域latent，则按单独生成目标计算flow loss，不将两种loss混到同一输出。

对保持空间对齐的编辑，可以用由速度反推的干净估计：

\[
\hat z_0=z_t-tu_\theta,
\qquad
\mathcal{L}_{\mathrm{preserve}}=
\operatorname{mean}_{i\in\mathrm{keep}}
\lVert\hat z_{0,i}-z_{\mathrm{source},i}\rVert_1.
\]

初始保持权重0.1仅作为候选。它会鼓励复制，必须同时检查目标变化成功率；改变相机、姿态和全局光照的样本不使用不适用的逐像素保持目标。最终严格外区仍靠源像素合成，而不是把这项损失解释为零漂移保证。

\section{S3：区间、多镜头与三十秒联合生成}\label{seedance-m20-s3ux533aux95f4ux591aux955cux5934ux4e0eux4e09ux5341ux79d2ux8054ux5408ux751fux6210}

保持原联合去噪结构，按8、16、30秒课程扩展。长阶段可用30\%、30\%、40\%的时长桶比例，并在每个桶内继续混合参考、编辑与基础任务。训练分辨率先用合法较低桶，固定对照后再提高；不能同时改变所有维度而无法定位退化。

caption包含global角色/环境与shot局部事件，事件时间映射到token物理时间；源条件和目标的不同噪声状态通过token条件明确传入。联合模型先不引入因果cache；先证明长时间的身份、事件、道具和声画关系能保持。

原生训练sampler的长度shift会线性外推，长序列直接照抄可能过度集中在纯噪声附近。短片保留原生对照；长阶段候选取90\%有界shifted logit-normal、10\%均匀sigma，shift限制在0.95至2.05，并固定配套推理策略。计算依据与有效样本损失细节见 \hyperref[seedance-m25]{数据与训练规格}，这项新课程需要独立比较。

出口包括：长片后半段不集中崩坏；角色出场区间和声音归属正确；参考容量增加有实际收益；局部返工不破坏其它镜头。阈值由真实任务在训练前确定，本文不虚构通过率。

\section{S4：单独训练因果续写模型}\label{seedance-m20-s4ux5355ux72ecux8badux7ec3ux56e0ux679cux7eedux5199ux6a21ux578b}

从联合模型派生新版本，以固定时间块为生成单位，视频和音频使用同一块划分。历史为干净条件，当前块为活动噪声；块间self-attention与音视频cross-attention都不允许读取未来。

先真实历史，再加入受控扰动历史；自生成历史由固定teacher在相同历史与下一事件条件下重标后续，并过滤状态冲突，不能任意接上原片真实未来。其比例逐渐从30\%提高到70\%只是试验课程，历史梯度截断。若错误累积压过训练信号，应回到较短rollout和更可靠历史，而不是继续增大时长。

缓存由固定条件下的历史层状态建立，每块完成后按干净零噪声历史重算再提交。提示、参考、模型、位置或过去编辑改变后重建；当前块每步仍计算QKV。原decoder可能需要未来上下文，首版采用离线整段解码。只有训练/推理依赖一致并完成独立验证，才可把它称为可缓存AR模型。此前的重叠尾段条件续写保留为可对照基线。

\section{S5：偏好、蒸馏与推理优化}\label{seedance-m20-s5ux504fux597dux84b8ux998fux4e0eux63a8ux7406ux4f18ux5316}

先用人工确认的失败和偏好做高质量SFT，再校准奖励模型。优化时视频动作、身份、声音、时间控制和非目标保持分开打分，不能以单一美感奖励抵消硬约束失败。

只有任务分布稳定后才做少步数student。多模态CFG分别控制文本与跨模态影响：

\[
\hat u=u(c,m)+s_c[u(c,m)-u(\varnothing,m)]
+s_m[u(c,m)-u(c,\varnothing)].
\]

这里的去条件分支必须被训练支持，且三个分支都产生网络求值成本。先保留原模型原生guidance，再单独增加参考strength；不要第一版把全部guidance组合扩成大量无对照分支。

\section{资源估计与停止规则}\label{seedance-m20-ux8d44ux6e90ux4f30ux8ba1ux4e0eux505cux6b62ux89c4ux5219}

19B主干的全参AdamW训练，若按BF16权重/梯度、FP32主副本与两组矩估算，每参数约16字节，仅这些状态就约304GB，尚不含activation、VAE、文本与临时张量。分片可以降低每卡占用，却不消除总状态和通信。LoRA/冻结主干会改变这笔账，实际需求需单独估算。

每阶段先设计短pilot，再按真实step耗时和验证成本决定训练量。这套新组合尚无GPU速度、训练小时或质量实测，论文作者的算力数字也不能直接当成它的预算。

奖励升而人工下降、动作趋静、主体串用、音色混合、长片后段错误聚集，都构成停止扩大训练的理由。恢复点保存所有新增模块、LoRA、optimizer、scheduler、随机状态、样本游标、数据/任务混合/sigma版本，避免回滚后已无法解释比较结果。原trainer主要提供权重保存路径，完整恢复状态需要另行实现。

真实训练入口以 \href{https://github.com/Lightricks/LTX-2/tree/bd92a5f408eb8c99f8c05e9fe2fcaaec8b56b638/packages/ltx-trainer}{LTX trainer} 为起点；自定义参考/源条件和因果化需要新增策略。它提供实施基础，不表示当前仓库已包含本文所有任务。
